\section{The bit decoder and a separating functional}\label{sec:decoder}

Here \(m\ge0\) is arbitrary and \(n=2^\ell\). We impose \(\ell\ge2\)
only for the old-base transfer and separation theorem.
The old compact bit system is
\begin{equation}\label{eq:compact-base}
 \mathcal Q_{m,\ell}
  =\{E_{ii'}:i<i'\in[m]\}\cup\Bool_b,
 \qquad
 E_{ii'}=\prod_{t=1}^{\ell}(1-b_{it}-b_{i't}).
\end{equation}
The degree of each collision generator is \(\ell\).
Its vanishing asserts that the two bit rows encode different labels.

Identify unary columns with labels \(z\in\{0,1\}^{\ell}\), and define the
ordinary ring map
\begin{equation}\label{eq:decoder}
 \tau(b_{it})=\sum_{z:z_t=1}X_{iz}.
\end{equation}

\begin{lemma}[Two-row interpolation]\label{lem:interpolation}\lean{claims/TwoRowInterpolation.lean}
For any common row width \(N\ge0\), let \(\mathcal U_{i,i'}\) consist of Booleanity, same-row exclusions,
and the two row-sum equations for distinct rows \(i,i'\).
For a polynomial \(R\) supported on these rows and of degree at most \(d\),
\[
 R-\sum_{z,w}R(e_z,e_w)X_{iz}X_{i'w}
     \in\Ideal_{\max\{d,2\}}(\mathcal U_{i,i'}).
\]
Here \(e_z,e_w\) are the corresponding one-hot assignments.
Column exclusions are not used in this interpolation.
\end{lemma}
\begin{proof}
Reduce powers and same-row collisions. The remainder is a linear
combination of \(1\), individual cells of the two rows, and products
containing one cell from each row. This reduction costs at most \(d\).
Replace the constant by \(\rho_i\rho_{i'}\), using
\[
 1-\rho_i\rho_{i'}=(1-\rho_i)+\rho_i(1-\rho_{i'}),
\]
and replace each remaining single cell by its product with the other
row sum. These changes have certificates of degree at most two.
The result is bilinear, and its coefficient at \(X_{iz}X_{i'w}\)
is its value at \((e_z,e_w)\), namely \(R(e_z,e_w)\).
\end{proof}


\begin{lemma}[A polynomial left inverse for the decoder]\label{lem:decoder-injective}
\lean{claims/CompactBitDecoder.lean}
For every \(m,\ell\ge0\), the decoder \(\tau\) is injective and preserves
ordinary total degree.
\end{lemma}
\begin{proof}
Define a linear substitution \(\lambda\) on unary variables by sending
\(X_{i,e_t}\) to \(b_{it}\), where \(e_t\) is the unit bit label, and
all other unary variables to zero. The unit labels are distinct and
\(\lambda\tau(b_{it})=b_{it}\), so \(\lambda\tau\) is the identity
on polynomials. Both substitutions have degree at most one; applying the
degree inequality in both directions proves equality, as well as injectivity.
\end{proof}

\begin{lemma}[Degree-preserving old-base transfer]\label{lem:decoder}
\lean{claims/CompactBitDecoderTransfer.lean}
For \(m\ge0\), \(\ell\ge2\), and every \(B\ge0\),
\[
 q\in\Cons_B(\mathcal Q_{m,\ell})
 \quad\Longrightarrow\quad
 \tau q\in\Cons_B(\mathcal F^{\mathrm{fun}}_{m,n}).
\]
\end{lemma}
\begin{proof}
Bit Booleanity has a degree-two image certificate:
\[
 \tau(b_{it}^2-b_{it})=\sum_{z:z_t=1}(X_{iz}^2-X_{iz}).
\]
At a two-row one-hot assignment \((e_z,e_w)\), \(\tau(E_{ii'})\)
is one exactly when \(z=w\). Lemma~\ref{lem:interpolation} therefore gives
\[
 \tau(E_{ii'})-\sum_zX_{iz}X_{i'z}
       \in\Ideal_\ell(\mathcal U_{i,i'}).
\]
The remaining terms are column exclusions of degree two. Thus every axiom
image has a certificate through its original degree. Replay the source
proof using Lemma~\ref{lem:substitution} with \(T=1\), \(w=1\).
Only axioms actually admitted by the source degree ceiling must be replayed;
there is no additional assumption \(B\ge\ell\).
\end{proof}


Let \(\mathcal L_k\) be the space spanned by monomials of degree at most
\(k\) containing at most one bit variable from each row.

\begin{lemma}[Row-linear dimension]\label{lem:row-dimension}
\lean{claims/RowLinearPolynomialSpace.lean}
Over any field, for all \(m,\ell,k\ge0\),
\begin{equation}\label{eq:row-linear-dimension}
 \dim\mathcal L_k=\sum_{j=0}^k\binom mj\ell^j.
\end{equation}
Every nonzero polynomial in this space has a nonzero monomial coefficient
of degree equal to its total degree.
\end{lemma}
\begin{proof}
Ordinary monomials form a basis. To choose a row-linear monomial of degree
\(j\), choose its \(j\) rows and one of \(\ell\) coordinates in each.
Terms with \(j>m\) contribute zero. The last assertion follows by choosing
a maximal-degree element of the finite nonempty support.
\end{proof}

\begin{lemma}[Binary cube degree drop]\label{lem:cube-drop}
\lean{claims/BinaryCubeDegreeDrop.lean}
Partition some source variables into \(t\) indexed groups. For each
\(\epsilon\in\F^t\), substitute scalars depending only on \(\epsilon_i\)
for variables in group \(i\), and fixed affine target polynomials for
all remaining variables. Write \(\phi_\epsilon\) for this substitution
and \(\Delta p=\sum_\epsilon\phi_\epsilon(p)\). Then
\[
 \deg\Delta p\le\max\{\deg p-t,0\},\qquad
 \deg p<t\Longrightarrow\Delta p=0.
\]
If every source generator \(F\) has a cube-independent image \(r_F\)
which is zero or a target generator, then
\[
 \Delta\Ideal_B(\mathcal F)\subseteq
       \Ideal_{\max\{B-t,0\}}(\mathcal G).
\]
The variable sets may be arbitrary; only the cube index set must be finite.
\end{lemma}
\begin{proof}
A monomial missing one group has equal evaluations in pairs obtained by
flipping that cube coordinate, so its sum is zero. A surviving monomial
uses at least \(t\) selected variables, which become scalars; the remaining
factors have total degree at most its original degree minus \(t\).
For a legal multiple \(qF\), cube independence gives
\(\Delta(qF)=(\Delta q)r_F\). If \(r_F=0\) it vanishes; otherwise it is
one target generator multiple, whose degree is at most \(B-t\) when
\(B\ge t\). When \(B<t\), it is zero. Extend by linearity.
\end{proof}

\begin{lemma}[Row-cube restriction and coefficient isolation]\label{lem:row-cube}
\lean{claims/DisjointCoordinatePairs.lean}
\lean{claims/RowCubeRestriction.lean}
\lean{claims/RowCubeCoefficientIsolation.lean}
Choose \(t\) distinct rows and, for each, two labels differing only in
one prescribed bit. If these pairs are disjoint, the one-hot substitutions
on the selected rows, together with deletion of their \(2t\) columns on
other rows, give
\[
 \Delta\Ideal_B(\mathcal F^{\mathrm{fun}}_{m,n})
 \subseteq\Ideal_{\max\{B-t,0\}}
       (\mathcal F^{\mathrm{fun}}_{m-t,n-2t}).
\]
For a row-linear \(f\) of degree at most \(t\), this cube applied to
\(\tau f\) is the constant coefficient of the monomial selecting those
rows and prescribed bits. Such disjoint pairs can be chosen whenever
\(t\ge1\) and \(4(t-1)<2^\ell\).
More generally, the coefficient-isolation identity holds for arbitrary
cube-independent polynomial images of the unselected bits, provided each
selected bit has its prescribed base value plus its selected cube coordinate,
and the other bits in that row stay constant.
\end{lemma}
\begin{proof}
On a selected row, each substitution is a one-hot assignment; on every
other row, deleted columns become zero and surviving variables are relabeled.
Booleanity, row exclusions and column exclusions consequently map to zero
or the corresponding residual generator. Disjointness prevents a collision
between selected rows or between selected and residual rows. A selected
row-sum equation maps to zero, and every other row-sum equation maps to
its residual counterpart. Lemma~\ref{lem:cube-drop} applies.

A row-linear monomial missing a selected row cancels by a cube flip. If it
uses all \(t\) selected rows, the degree bound forces it to use exactly
one variable in each and none outside. The two-value sum in each row is
one for the prescribed bit and zero for every other bit. This isolates
exactly the claimed coefficient, independently of the images outside the
selected rows. Finally, after selecting \(j\) pairs, their \(2j\) vertices
forbid at most \(2j\) of the \(2^{\ell-1}\) pairs in any specified direction.
The strict packing bound permits the next choice.
\end{proof}

\begin{theorem}[Row-linear separation through a larger PC degree]
\label{thm:row-separation}\lean{claims/CubeResidualDualSeparation.lean}
Let \(m\ge0\), \(\ell\ge2\), \(n=2^\ell\), and suppose
\(1\le k\le B\), \(4(k-1)<n\), and \(2B-1\le n\).
For every nonzero \(f\in\mathcal L_k\), there is a linear functional
\(\mu\) on unary polynomials through degree \(B\) such that
\[
 \mu(\Cons_B(\mathcal F^{\mathrm{fun}}_{m,n}))=0,
 \qquad \mu(\tau f)=1.
\]
It can be chosen with \(\mu(1)=0\) if \(\deg f>0\), and
\(\mu(1)=1\) if \(\deg f=0\). Consequently
\[
 \tau f\notin\Cons_B(\mathcal F^{\mathrm{fun}}_{m,n}),\qquad
 f\notin\Cons_B(\mathcal Q_{m,\ell}).
\]
\end{theorem}
\begin{proof}
A nonzero constant over \(\F\) is one, so that case follows from a
normalized matching functional. Otherwise put \(t=\deg f\) and choose a
nonzero degree-\(t\) coefficient. Lemma~\ref{lem:row-cube} supplies its
coordinate pairs and residual operator \(\Delta\), with \(\Delta\tau f=1\).
Because \(n\) is divisible by four, \(4(k-1)<n\) implies \(k\le n/4\).
Thus \(n-2t\ge1\), and
\[
 n-2t\ge2(B-t)-1.
\]
Theorem~\ref{thm:moments} supplies a normalized residual annihilator
\(\lambda\) through \(B-t\), even if there are fewer than \(B-t\)
remaining rows. Set \(\mu=\lambda\circ\Delta\).
The residual NS transport and Corollary~\ref{cor:stability} show that
\(\mu\) annihilates the old unary PC space. Coefficient isolation gives
\(\mu(\tau f)=1\), whereas \(\Delta1=0\) gives \(\mu(1)=0\).
Finally apply Lemma~\ref{lem:decoder} for the bit consequence.
\end{proof}
